\begin{abstract}
Additive coagulation with constant per-particle fragmentation can conserve
particle count while the number and mass distributions move toward
opposite size boundaries. We prove fractional-moment inequalities that
quantify this separation, with coefficients that are sharp as uniform
instantaneous constants, for measurable time controls and
parent-dependent expected daughter measures. These coefficients are
Lyapunov bounds, not asymptotic rates: the explicit Golovin--Borel
solution decays several times faster. An auxiliary process with
the normalized number marginals accumulates only a finite upward
log-size correction and undergoes finitely many coagulations almost
surely, even when its expected total number of coagulations is infinite.
For parent-independent daughter fractions, this gives an exact finite
transport cost to a fragmentation reference and transfers compound-Poisson
scaling limits. We bound the last auxiliary coagulation time, prove daughter-law
envelopes that are extremal over all mean-conserving daughter laws, and
obtain exact representations of critical last-event observables in terms
of the unknown law. The open critical exponent reduces to the decay rate
of the deterministic number--mass overlap, and archived simulations
suggest it is near half the coagulation rate. At critical
balance $\sigma=\lambda m$, a physical finite population's normalized
count remains accurate on sublinear time
horizons although mass-distribution error approaches its maximum by
logarithmic times from identical initial data; this is a corollary of
the moment bound and the material ceiling of a finite vessel. Finally, a
Fourier factorization with an exponentially small remainder in time
yields structural fragmentation identification from ideal late-time
distributions on a low-frequency window, at a certified rate whose
asymptotic optimality is not asserted. A separate independent sampling model gives a pointwise
fixed-frequency ratio bound and a sufficient observation-time balance. These results distinguish bulk observability,
auxiliary path asymptotics, and finite-population validity.
\end{abstract}

\section{Introduction}
\label{sec:introduction}

Count, total mass, and mean size do not determine the evolution of a
coagulating and fragmenting population. For the additive kernel
$K(x,y)=\lambda(x+y)$ and constant per-particle fragmentation rate
$\sigma$, mass conservation gives
\[
 N'(t)=(\sigma-\lambda m)N(t).
\]
At the critical balance $\sigma=\lambda m>0$, all three bulk observables
remain fixed. Nevertheless, almost all particles by number become small,
while almost all material is carried by increasingly large particles.
The distinction between selecting a particle and selecting a unit of
material is therefore central to both asymptotic analysis and observation.
Throughout, ``constant fragmentation rate'' means constant per-particle
selection rate; a constant binary fragmentation density can instead
have a size-dependent integrated selection rate.

Population-balance modeling and the recovery of aggregation and breakage
rates from dynamic size distributions have a long history; see
\citet[Chapters 2, 6, and 7]{Ramkrishna2000} for the modeling, inverse,
and stochastic viewpoints, and \citet{Aldous1999,Wattis2006} for
mean-field surveys. Pure additive coagulation is explicitly solvable
\citep{Golovin1963}, has a probabilistic interpretation through
branching structures and the additive coalescent
\citep{AldousPitman1998,DeaconuTanre2000,Bertoin2002}, and a complete
scaling theory \citep{MenonPego2008}. Fractional-moment estimates are
standard tools in the analysis of Smoluchowski equations
\citep{EscobedoMischlerPerthame2002,FournierLaurencot2005,FournierLaurencot2006,EscobedoMischlerRodriguezRicard2005}. Constant count in an evolving population
with aggregation and fragmentation is also established.
\citet{PatilAndrews1998} and \citet{Lage2002} provide constant-count
examples, as explicitly discussed in the introduction of
\citet{McCoyMadras2003}, whose solution allows count to vary.
The results here concern quantitative separation of sampling laws,
its pathwise consequences for the additive model, and the implications
for finite populations and inverse observations.

\subsection{Results and their relation to established ingredients}

The main estimate is a decay inequality for every normalized
fractional moment
$\mathcal A_p=M_p/(N^{1-p}m^p)$, $0<p<1$.
These quantities compare the number law $\eta_t=n_t/N(t)$ with the mass
law $\pi_t=xn_t/m$; at $p=1/2$ they give their Hellinger affinity.
\Cref{thm:affinity} proves
\[
 \mathcal A_p(t)\leq\mathcal A_p(0)
 \exp\!\left[-a_p\int_0^t\lambda(s)m\dd s
             -a_{1-p}\int_0^t\sigma(s)\dd s\right],
 \qquad a_p=1+p-2^p.
\]
The coefficients are sharp as uniform instantaneous coefficients over
the class of solutions and daughter laws, in the explicit sense stated in
the theorem. They are not asymptotic rates: for the Golovin--Borel
solution the exact exponent of $\mathcal A_{1/2}$ is about six times
$a_{1/2}$, and the critical simulations of \cref{sec:finite} show a
similar gap. The pair inequality behind the estimate is a weighted
variant of the subadditivity-defect estimates used in the works cited
above; the contribution is the explicit constant with its equality case
and its uniform use for parent-dependent daughters. At criticality the estimate
proves number concentration at zero, mass escape to infinity, and
nonstationarity without a logarithmic moment assumption. Mass remains
conserved at each finite time.

The same estimate controls the nonlinear contribution to log size.
Probabilistic representations of coagulation equations and changes
between sampling weights are established
\citep{DeaconuTanre2000,DeaconuFournierTanre2002}.
We construct directly an auxiliary number process with marginals
$\eta_t$, including measurable parent-dependent daughters.
Its accumulated upward log increment is finite in mean and almost
surely. A daughter-product martingale then proves that its total
number of coagulation jumps is finite almost surely, although the
mean can be infinite. With selfsimilar daughters, coupling to the
independent fragmentation reference gives an exact extended
Wasserstein cost. Classical compound-Poisson limit theory transfers
to the nonlinear number law, with the daughter log moments needed
for each limit stated separately. Pure-fragmentation logarithmic
asymptotics already appear in
\citet{Bertoin2003,DoumicEscobedo2016}; the finite coagulation
correction supplies the transfer here.

The last auxiliary coagulation gives a stronger path-law comparison.
\Cref{sec:last-event} bounds its tail under arbitrary controls, proves
all-constant-rate daughter-law envelopes by a convexity argument, and
compares them with the number--mass overlap. At criticality the future
hazard is an exponential functional of a compound-Poisson subordinator
\citep{BertoinYor2005}, and the equal-split formula is the Poissonian
exponential functional of \citet{BertoinBianeYor2004}. We derive exact
representations of last-event observables and a density in terms of
the unknown law, together with upper and lower calendar-time tail
bounds. The exponent between these bounds is not determined; the
envelopes reduce it to the decay rate of the deterministic overlap, and
the archived simulations suggest the value $b/2$ of the pure-coagulation
benchmark. The Golovin--Borel solution \citep{Golovin1963}, recorded by
\citet{Bertoin2009}, provides that benchmark. These established
formulas are used with explicit attribution.

A finite physical population must obey a total-material ceiling that
the continuum equation does not share. Differences between finite
stochastic coagulation--fragmentation systems and mean-field predictions
are already documented, for example by
\citet{DOrsognaLeiChou2015} in a finite-cluster assembly model.
For the present critical additive model, \cref{sec:finite} records an
explicit distinction in growing time: starting from exactly the same
empirical initial measure, mass cumulative distributions have nearly
maximal discrepancy at least by times proportional to the logarithm of
population scale, while normalized count is uniformly accurate on every
sublinear horizon. The discrepancy statement is a corollary of the
moment bound and the material ceiling of a finite vessel; only the count
statement uses the particle dynamics. Nine reproducible particle runs
illustrate separation diagnostics and the gap between certified and
observed decay rates. They do not compute a continuum comparison.

Finally, the integrable nonlinear log correction gives a Fourier
factorization whose remainder is exponentially small in time, at the
certified rate $\kappa(b+\sigma)$. Local late-time quotients on a
low-frequency window identify the fragmentation exponent and, by
one-sided transform uniqueness, the expected daughter fraction measure.
The second step is analytic continuation from a short interval and is
severely ill-posed. This sits within established
inverse population-balance work: \citet{DoumicEscobedoTournus2018}
uses asymptotic fragmentation profiles,
\citet{DoumicEscobedoTournus2024} uses short-time distributions,
and \citet{MirzaevByrneBortz2016} estimates daughter distributions in
an aggregation--fragmentation model. Two-time quotient cancellation
and distinguished logarithms have a direct precedent in
\citet[Section 2.2]{Garnier2024}. The contribution of
\cref{sec:fourier} is the proved nonlinear remainder and its
identification consequence for additive coagulation. The statistical
application is a pointwise bound under independent sampling,
with no claim of stable recovery of a whole daughter measure or
of statistical optimality.

\subsection{Scope, assumptions, and organization}

Three different probabilistic objects occur. The auxiliary process
represents one-time continuum number laws; its jumps are not a physical
material genealogy. The finite particle system describes an actual
mass-conserving vessel with interacting particles. The empirical Fourier
sampling bound assumes independent draws from deterministic continuum
marginals. No independence or approximation between these objects is
implicit in the results.

\Cref{tab:assumptions} records the main changes of assumptions.
All deterministic estimates are conditional on the mass-conserving
weak-solution framework in \cref{sec:model}.
\Cref{thm:existence,app:existence} prove existence and uniqueness with
finite initial second moment, while most asymptotic estimates use only
finite number and mass once a solution is supplied.
A finite second size moment does not imply an integrable negative
log size. No logarithmic moment is silently inherited from that
existence theorem.

\begin{table}[htbp]
\centering\small
\begin{tabular}{p{0.29\textwidth}p{0.64\textwidth}}
\toprule
Results & Additional scope or assumptions \\
\midrule
Sampling separation and finite auxiliary correction
& Locally integrable time controls; measurable parent-dependent expected
  daughters with count two and conserved mean mass. \\
Transport and logarithmic limits
& Parent-independent fractions for the independent reference; fixed rates
  and fraction law for stationary increments; daughter log moments as stated. \\
Daughter envelopes and critical representations
& Constant rates; critical formulas require $\sigma=\lambda m$;
  equal splitting only for the formulas labeled as such. \\
Finite-population comparison
& Critical rates and eventwise complementary positive binary splits;
  continuum and particles start from the identical empirical measure. \\
Fourier identification and sampling
& Constant $\lambda>0$, $\sigma\geq0$, and fixed parent-independent $B$;
  exact distributions for identification, independent samples for the sampling bound. \\
\bottomrule
\end{tabular}
\caption{Assumption roadmap. Individual statements specify all moment,
observation, and solution hypotheses.}
\label{tab:assumptions}
\end{table}

\Cref{sec:boundaries} adds the complete fixed-mass count-neutral
classification, rigidity from one further moment, and the failure of
fractional-moment monotonicity beyond the additive kernel. These results
clarify which conclusions depend on the additive structure.
Elementary finite-grid preparation and initial-rate algebra, including
dilution and source ambiguities, are collected in a supplementary note
distributed with the numerical supplement.
The appendices contain forward well-posedness, product error
certification, finite-count formulas, and unbounded-test justifications.
\Cref{sec:discussion} separates the proved conclusions from the remaining
asymptotic and statistical questions.
